\cotefile{47}
\scanpage{2}
Produit de convolution sur les VA et transformation de Fourier \unsure{sur|dans} les VA

% archive sheet number: 1
\scanpage{3}
\uline{Proposition} \unsure{1|!} Soient $S$ un schéma, $A$ un schéma abélien sur $S$, $X$ un schéma sur $S$, \deleted{\unsure{tel que pour toute|tel que pour tout} \unsure{fibre géométrique}} \inserted{$L$ un faisceau inversible sur} $A\times_S X$, \deleted{Considérons l'homomorphisme} trivial sur $e_A\times_S X$, et tel que \deleted{\unsure{pour que tout|pour tout} \unsure{$L_x$} \illegible{1}} pour tt $x\in X$, le faisceau $L_x$ sur $A_x=A_s\otimes_{k(s)}k(x)$ ($s$ le pt de $S$ sous $x$) soit alg. équivalent à zéro. Considérons l'application
% slanted
\margin{peut se généraliser au cas $A$ un schéma en groupes \inserted{\unsure{plat|plats}}, $L$ sur $A\times_S X$, satisfaisant $(\pi_A\times\mathrm{id}_X)^*(L)\simeq\mathrm{pr}_{13}^*(L)\otimes\mathrm{pr}_{23}^*(L)$ et $F$ et $G$ sur $A$ à supports propres sur $S$ et \deleted{\unsure{parfait}} \inserted{$F,G$} $\in D^{b}_{\mathrm{qcoh}}(A)$\ldots}
\[
\varphi_L\colon K^{\cdot}(A)\longrightarrow K^{\cdot}(X)
\]
définie par \deleted{l'application additive} le foncteur de catégories triangulées \deleted{$\unsure{\varphi}$} $\Phi\colon D_{\mathrm{parf}}(A)\unsure{\mathrel{-}|\longrightarrow}D_{\mathrm{parf}}(X)$
\[
\Phi_L(F)=\mathbb{R}\mathrm{pr}_{2*}\bigl(\mathbb{L}\mathrm{pr}_1^*(F)\otimes L\bigr).
\]
Alors $\varphi_L$ transforme produit de Pontryagin $*$ en produit ordinaire. Plus précisément, on a un isomorphisme de foncteurs de catégories triangulées
\[
\Phi_L(F\overset{\mathbb L}{*}G)\simeq\Phi_L(F)\overset{\mathbb L}{\otimes}\Phi_L(G)
\]
$(F,G\in\mathrm{Ob}\,D_{\mathrm{parf}}(A))$.
% slanted
\margin{on peut enfin généraliser en un énoncé concernant \deleted{\unsure{un|des} \unsure{couple|complexes}} de $C$-torseurs \diagram{\begin{tikzcd}[ampersand replacement=\&,column sep=tiny,row sep=tiny]P\times Q\arrow[r] \& R\\ C\times C\arrow[r] \& C\end{tikzcd}}, et \unsure{même} on peut laisser tomber $C$, et partir de faisceaux $L,L',L''$ sur $R\times X,P\times X,Q\times X$, et $\pi\colon P\times Q\unsure{\mathrel{-}|\to} R$, \deleted{\unsure{tels que}} tels que $\deleted{\unsure{R^{*}}}(\pi\unsure{\times|\mathbin{\&}|\alpha}\mathrm{id}_X)^*(L)\simeq\mathrm{pr}_{13}^*(L')\,\mathrm{pr}_{\unsure{22|23}}^*(L'')$}

Si le th. est démontré dans la situation donnée, il en résulte qu'il est vrai en tte situation \inserted{$(A,X',L')$} qui s'en déduit par image inverse, à l'aide d'un $S$-morphisme $f\colon X'\unsure{\mathrel{-}|\to}X$, car $\Phi_{L'}\simeq\mathbb{L}f^*(\Phi_L)$, et $\mathbb{L}f^*$ commute à $\overset{\mathbb L}{\otimes}$.

Soit $B$ le schéma abélien dual de $A$, et $L_0$ le faisceau de Weil sur $A\times_S B$. Alors l'hypothèse sur $L$ implique que $L$ est $\simeq$ l'image inverse de $L_0$ par un $S$-morphisme $X\longrightarrow B$. Donc on est ramené au cas où $X=B$, $L=L_0$. \deleted{\unsure{Or|On}} On a des isom. fonctoriels
\[
F\overset{\mathbb L}{\deleted{\otimes}{*}}G=\mathbb R\pi_{\unsure{*|\alpha}}\bigl(F\overset{\mathbb L}{\otimes}_S G\bigr)\qquad\text{et}\tag*{$(*)$}
\]
% archive sheet number: 2
\scanpage{4}
\diagram{\begin{tikzcd}[ampersand replacement=\&,column sep=small,row sep=small]
A\times A \arrow[d,"\pi"] \arrow[dr,phantom,"{\text{cart}}"{description,draw,ellipse,inner sep=1pt}] \& A\times_S A\times_S B \arrow[l,"\mathrm{pr}_{12}"'] \arrow[d,"{\pi\times \mathrm{id}_B}"] \arrow[dr,"\mathrm{pr}_3"] \& {}\\
A \& A\times_S B \arrow[l,"\mathrm{pr}_1"] \arrow[r,-,"\mathrm{pr}_2"'] \& B
\end{tikzcd}}

\diagram{\begin{tikzcd}[ampersand replacement=\&,column sep=small,row sep=small]
{} \& A\times_S A\times_S B \arrow[r,"\mathrm{pr}_3"] \arrow[d,"{\mathrm{id}_{A\times A}\times\Delta_B=f}"'] \arrow[ddl,bend right=39,"\mathrm{pr}_{12}"'] \arrow[dr,phantom,"{\text{cart}}"{description,draw,ellipse,inner sep=1pt}] \& B \arrow[d,"\Delta_B"]\\
{} \& A\times_S A\times_S B\times_S B \arrow[r,"\mathrm{pr}_{34}"'] \arrow[dl,"\mathrm{pr}'_{12}"] \& B\times_S B\\
A\times_S A \& {} \& {}
\end{tikzcd}}

\[
\begin{aligned}
\Phi(F\overset{\mathbb L}{*} G)&\overset{\text{déf}}{\simeq}\deleted{\unsure{\mathbb{R} \mathrm{pr}_{2*}(\mathbb{L} \mathrm{pr}_1^*(F\overset{\mathbb L}{*} G)\otimes L)}}\inserted{\mathbb{R} \mathrm{pr}_{2*}(\mathbb{L} \mathrm{pr}_1^*(F\overset{\mathbb L}{*} G)\otimes L)}.\\
&\overset{(*)}{=}\mathbb{R} \mathrm{pr}_{2*}\bigl(\mathbb{L} \mathrm{pr}_1^*(\pi_*(F\overset{\mathbb L}{\otimes} G))\inserted{\otimes L}\bigr)\\
&\overset{\text{chgt. de base}}{\simeq}\mathbb{R} \mathrm{pr}_{2*}\bigl(R(\pi\times \mathrm{id}_B)_*(\mathbb{L} \mathrm{pr}_{12}^*(F\overset{\mathbb L}{\otimes}_S G))\otimes L\bigr)\\
&\overset{\text{formule de proj.}}{\simeq}\mathbb{R} \mathrm{pr}_{2*}\bigl(\mathbb{R}(\pi\times \mathrm{id}_B)_*(\mathbb{L} \mathrm{pr}_{12}^*(F\overset{\mathbb L}{\otimes}_S G)\otimes(\pi\times \mathrm{id}_B)^*(L))\bigr)\\
&\overset{\text{transit.}}{\simeq}R \mathrm{pr}_{3*}\bigl(\mathbb{L} \mathrm{pr}_{12}^*(F\overset{\mathbb L}{\otimes}_S G)\otimes(\pi\times \mathrm{id}_B)^*(L)\bigr)
\end{aligned}
\]

Or on a (th. du carré) un isom can.
\[
(\pi\times \mathrm{id}_B)^*(L)\simeq \mathrm{pr}_{13}^*(L)\,\mathrm{pr}_{23}^*(L)
\]
où $\mathrm{pr}_{13},\mathrm{pr}_{23}\colon A\times A\times B\unsure{\mathrel{-}|\to} A\times B$ sont respectivmt $\mathrm{pr}_1\times \mathrm{id}_B$ et $\mathrm{pr}_2\times \mathrm{id}_B$.

On en conclut
\[
\mathbb{L} \mathrm{pr}_{12}^*(F\overset{\mathbb L}{\otimes}_S G)\otimes(\pi\times \mathrm{id}_B)^*(L)\simeq \mathbb{L} f^*\bigl(\underline{\tikz[remember picture,baseline]{\coordinate (quatretermes) at (0,0);}\mathrm{pr}_{12}^{\prime *}(F\overset{\mathbb L}{\otimes}_S G)\otimes \mathrm{pr}_{13}^{\prime *}(L)\otimes \mathrm{pr}_{24}^{\prime *}(L)}\bigr)
\]
où $\mathrm{pr}'_{13},\mathrm{pr}'_{24}$ sont les deux projections $A\times_S A\times_S B\times_S B\longrightarrow A\times_S B$. On a d'ailleurs
\[
\mathrm{pr}_{12}^{\prime *}(F\overset{\mathbb L}{\otimes}_S G)\otimes \mathrm{pr}_{13}^{\prime *}(L)\otimes \mathrm{pr}_{24}^{\prime *}(L)
\overset{\text{ass. de }\unsure{\overset{\mathbb L}{\otimes}_S|\otimes_S|\otimes}}{\simeq}
\tikz[remember picture,baseline]{\coordinate (quatrehaut) at (0,0);}\,\mathrm{pr}_{13}^{\prime *}(F\otimes L)\overset{\mathbb L}{\otimes} \mathrm{pr}_{24}^{\prime *}(G\otimes L)
\]
\deleted{conclut} et on conclut
\[
\begin{aligned}
\Phi(F\overset{\mathbb L}{*} G)&\simeq R \mathrm{pr}_{3*}\bigl(\mathbb{L} f^*(\tikz[remember picture,baseline]{\coordinate (quatrebas) at (0,0);}\hspace{1em})\bigr)\\
&\overset{\text{chgt. de base}}{\simeq}\mathbb{L}\Delta_B^*\bigl(\mathbb{R} \mathrm{pr}_{34*}(\mathrm{pr}_{13}^*(F\otimes_S L)\overset{\mathbb L}{\otimes} \mathrm{pr}_{24}^{\prime *}(G\otimes_S L))\bigr)\\
&\overset{\text{Künneth}}{\simeq}\mathbb{L}\Delta_B^*\bigl(\mathbb{R} \mathrm{pr}_2(F\otimes L)\overset{\mathbb L}{\otimes}_S R \mathrm{pr}_{2*}(G\otimes L)\bigr)\\
&\overset{\text{déf}}{\simeq}\mathbb{L}\Delta_B^*\bigl(\Phi(\unsure{L|F})\overset{\mathbb L}{\otimes}_S\Phi(G)\bigr)\\
&\simeq\Phi(F)\overset{\mathbb L}{\otimes}\Phi(G)
\end{aligned}
\]
\tikz[remember picture,overlay]{\draw (quatrebas) |- (quatrehaut); \draw[->] (quatrehaut) to[bend right=12] (quatretermes);}
cqfd.

\deleted{\uline{Proposition 2}} \inserted{\uline{Théorème}} Soient $A$, \deleted{deux} \inserted{$B$} schémas abéliens sur $S$\inserted{, de dim relative $n$}, \deleted{$\unsure{B|L}$} $L$ un Module inversible birigidifié sur $A\times B$, définissant une \uline{dualité} entre $A$ et $B$. \deleted{Considérons} \inserted{Alors le composé} des deux
\diagram{\begin{tikzcd}[ampersand replacement=\&]
K(A)\arrow[r,"{\unsure{\varphi|\Phi}_L}"] \& K(B)\arrow[r,"{\varphi_{{}^{s}L}}"] \& K(A)
\end{tikzcd}}
\deleted{\unsure{sont|est} \unsure{l'identité}} \inserted{$(-1)^n s_A^*$ (où $s_A$ est la symétrie de $A$)} (donc, par symétrie, on voit que $\Phi_L$ et $\inserted{(-1)^n}s_A^*\Phi_{{}^{s}L}$ sont des isomorphismes inverses l'un de l'autre).
% archive sheet number: 3
\scanpage{5}
Plus précisément, le foncteur composé
\diagram{\begin{tikzcd}[ampersand replacement=\&]
D_{\mathrm{parf}}(A)\arrow[r,"\Phi_L"] \& D_{\mathrm{parf}}(B)\arrow[r,"{\Phi_{{}^{s}L}}"] \& D_{\mathrm{parf}}(A)
\end{tikzcd}}
est \inserted{can.} isomorphe au foncteur \deleted{identique, \unsure{par suite}} \inserted{$\mathbb{L} s_A^*\deleted{\unsure{Rs}}\simeq R s_{A*}$ suivi de $\tau_A[-n]$}, \deleted{\unsure{ce qui|de sorte que|respectivement}} où $\tau_A\simeq\unsure{\bigwedge^{n}|\bar{\Lambda}|\hat{\Lambda}}t_A\simeq\bigwedge^{n}t_B$. Donc $\Phi_L$ et $\mathbb{L} s_A^*\Phi_{{}^{s}L}\inserted{=\mathbb{R} s_{A*}\Phi_{{}^{s}L}}$ sont des équivalences des catégories triangulées de $D_{\mathrm{parf}}(A)$ et $D_{\mathrm{parf}}(B)$.
% slanted
\margin{inutile de se borner aux complexes parfaits, on peut les prendre p.ex. dans $D^b_{\mathrm{qcoh}}\deleted{\unsure{A}}$}
\margin{\deleted{\unsure{des catégories|identique|identité}}}

\uline{Corollaire} \unsure{3|2} $\Phi_L\circ\omega_A[n]$ : \deleted{transforme} \inserted{$F\longmapsto\Phi_L(\unsure{F|\omega_A}\otimes\omega_A[n])\simeq\Phi_L(F)\otimes\omega_A[n]$ transforme} $\overset{\mathbb L}{\otimes}$ en $\overset{\mathbb L}{*}$ : \deleted{(et \unsure{inversement} $\overset{\mathbb L}{*}$ en $\overset{\mathbb L}{\otimes}$)}. D'où $=$ \deleted{$\unsure{(-1)^n\varphi_L}$} \deleted{$\unsure{x\mapsto\varphi_L(x}$} \inserted{$x\longmapsto(-1)^n\mathrm{cl}\,\omega_A\,\varphi_L(x)$} transforme produit ordinaire en \deleted{\unsure{cup}} produit de convolution.
\[
\begin{aligned}
\Phi_L(F\unsure{\overset{\mathbb L}{\otimes}|\otimes}G)\otimes\omega_A[n]&\simeq[\Phi_L(F)\otimes\omega_A[n]]*[\Phi_L(G)\otimes\omega_A[n])\\
\text{i.e.}\quad \Phi_L(F\overset{\mathbb L}{\otimes} G)&\simeq[\Phi_L(F)*\Phi_L(G)]\unsure{\otimes|\mathbin{\cdot}}\omega_A[n]
\end{aligned}
\]

Prouvons la prop. 2. Il suffit de définir un isom $\Phi_{{}^{s}L}\circ\Phi_L\simeq \mathrm{id}_{D_{\mathrm{parf}}(A)}$. Or le composé des foncteurs est de la forme $F\longmapsto \mathbb{R} \mathrm{pr}_{2*}(\mathrm{pr}_1^*(F)\overset{\mathbb L}{\otimes} M)$, où $M\in D_{\mathrm{parf}}(A)$ est donné par
\margin{$A\times B\times A$}
\[
M=\unsure{\mathbb{R}|\mathbb L R}\mathrm{pr}_{13*}\bigl(\mathrm{pr}_{12}^*(L)\,\mathrm{pr}_{23}^*({}^{s}L)\bigr)
\]
Il faut donc calculer ce $\mathbb{R} \mathrm{pr}_{13*}$, \deleted{et} Pour ceci, posons $C=A\times A$, et considérons l'isom $A\times B\times A\simeq\deleted{\unsure{C\times B\times}}C\times B$. Le module inversible $N=\mathrm{pr}_{12}^*(L)\,\mathrm{pr}_{23}^*({}^{s}L)$ sur $C\times B$ est \deleted{\unsure{muni de}} birigidifié de façon canonique, et correspond à un homomorphisme $C\xrightarrow{\psi}B^*$, qui n'est autre que l'hom. de $A\times A$ dans $B^*$ dont les \inserted{deux} composantes sont égales à $\varphi\colon A\unsure{\mathrel{-}|\to}B$ (l'hom de $A$ dans $B^*$ défini par $L$).
% archive sheet number: 4
\scanpage{6}
\diagram{\begin{tikzcd}[ampersand replacement=\&]
C\times_S B\arrow[r,"{\psi\times \mathrm{id}_B}"]\arrow[d,"\mathrm{pr}_1"']\arrow[dr,phantom,"{\text{cart}}"{description,draw,ellipse,inner sep=1pt}] \& B^*\times_S B^{\deleted{*}}\arrow[d,"\mathrm{pr}'_1"]\\
C\arrow[r,"\psi"'] \& B^*
\end{tikzcd}}
Ceci dit, on aura
\[
\deleted{R \mathrm{pr}_{2*}(N)}\qquad N\simeq(\psi\times \mathrm{id}_B)^*(W_B)
\]
($W_B$ le Module de Weil sur $B\times_S B^*$), d'où
\[
\begin{aligned}
\mathbb{R} \mathrm{pr}_{2*}(N)&\simeq \mathbb{R} \mathrm{pr}_{2*}(\psi\times \mathrm{id}_B)^*(W_B)\\
&\overset{\text{chgt de base}}{\simeq}\mathbb{L}\psi^*(\mathbb{R} \mathrm{pr}'_{2*}(W_B)).
\end{aligned}
\]

\uline{Lemme 1} Soit $B$ un schéma abélien sur $S$\inserted{, de dim relative $n$}, $W$ le faisceau de Weil sur $B^*\times_S B^{\deleted{\unsure{*}}}$, $p\colon B^*\times_S B^{\deleted{\unsure{*}}}\unsure{\mathrel{-}|\to}B^*$ la projection. Alors
\[
\mathbb{R} p^{\deleted{\unsure{*}}}_*(W)\simeq e_{\deleted{\unsure{B}}*}\bigl[\deleted{\unsure{\mathcal{O}\otimes\tau[-n]_{B^*}}}\tau_{B^*}[-n]\bigr)
\]
(isom. canonique), où $\tau_{B^*}=\bigwedge^n\underline{t}_{B^*}$ ($t_{B^*}$ le faisceau tangent à $B^*$ le long de la section $0$).
\margin{$\simeq e_*(\omega^{-1})$ où $\omega=\det^{*}(\underline{\Omega}^1_{B/\unsure{A|S}})$}

On en conclut un

\uline{Corollaire} Soient $C,B$ schémas abéliens sur $S$\inserted{, $B$ de dim relative $n$}, $N$ un Module inversible birigidifié sur $C\times_S B$, $\psi\colon C\longrightarrow B^*$ l'hom correspondant. Alors, on a
\[
\mathbb{R} \mathrm{pr}_{1*}(N)\simeq \mathbb{L}\psi^*(\unsure{\tau|e|e_*\tau}_{B^*}\deleted{\unsure{*}}[-n])\quad\deleted{\unsure{a}}
\]
En particulier, si $\psi$ est un épimorphisme, \deleted{\unsure{et si|donc}} donc plat, et si $K=\operatorname{Ker}\psi$, alors $K$ est un sous-schéma en groupe de $C$ plat de prés. finie sur $S$, donc \deleted{\unsure{définit un|défini} \unsure{complexe}} $\underline{O}_K$ est parfait sur $C$, et on a
\[
\mathbb{R} \mathrm{pr}_{1*}(N)\simeq\underline{O}_K[-n]\otimes_S\tau_{B^*}
\]

Dans le cas qui nous occupe, $C=A\times A$, \deleted{$\unsure{C^*\simeq A}$} $\varphi\colon A\simeq C^*$, et \deleted{\unsure{est}} par suite $K$ est l'\uline{antidiagonale} de $A\times_S A$, noyau de $\pi_A\colon A\times_S A\longrightarrow A$, qui est le graphe de la \uline{symétrie} $s_A$ de $A$. Donc \deleted{\unsure{la} \illegible{1}} \unsure{on gagne}, modulo dém. du Lemme 1.
% archive sheet number: 5
\scanpage{7}
\uline{\unsure{Démontrons le|Démonstration du} Lemme} 1 \quad Il sera une conséquence du

\uline{Lemme 2} \inserted{(Mumford)} Soit $\deleted{\unsure{A}}B$ schém. abélien sur un corps \unsure{[alg. clos ?]|(alg. clos ?)|(alg. clos)} $k$, et $L$ un Module inv. sur $B$ alg. équivalent à zéro, \deleted{\unsure{Al|Alors}} et non isomorphe à $\mathcal O_B$. Alors on a
\[
\deleted{\unsure{-}}\quad H^i(X,L)=0\quad\deleted{\unsure{\text{si }i>0}}\quad\text{pour \unsure{\deleted{tout}|tout} }i
\]
\[
\deleted{\text{et }\ H^0(X,L)=
\begin{cases}
0&\text{si }L\unsure{\not\simeq|\nsim}\underline O_B\\
k&\text{si }L\simeq\underline O_B
\end{cases}}
\]

\deleted{\uline{Démonstration du lemme 2} L'assertion sur le $H^0$ est triviale, et il reste à prouver que $H^i(X,L)=0$ si $i\unsure{>}$.}

Nous allons ici admettre le lemme 2, et allons montrer comment en \deleted{conclure} \inserted{déduire} le lemme 1. Nous voyons déja grâce à lemme 2 que $Rp_*(W)$ est concentré sur la section unité de $B^*$. Nous allons définir un homomorphisme canonique (où on pose $e_{B^*}=e$)
\[
Rp_*(W)\longrightarrow\deleted{\unsure{e_*(\underline O_S)[-n]}}e_*(\underline{\mathcal O}_S)\otimes_{\underline O_S}\tau_{B^*}[-n]
\]
ou, ce qui revient au $=$, \inserted{(comme $n$ est la dim relative de $B/S$, donc $R^ip_*(W)=0$ si $i>n$)} un homomorphisme canonique
\[
R^np_*(W)\longrightarrow\deleted{\unsure{e_*(\underline O_S)|\otimes_S}}e_*(\tau_{B^*})\deleted{\unsure{[-n]}}
\]
ou encore un hom.
\[
\deleted{e_{B}}\quad e^*(\mathbb{R}^n p_*(W))\longrightarrow\deleted{\unsure{\underline O_S|\otimes_S}}\tau_{B^*}
\]
\diagram{\begin{tikzcd}[ampersand replacement=\&]
B^*\times_S B\arrow[d,"p"] \& B\arrow[l,"{e'}"']\arrow[d,-,"q"]\\
B^* \& S\arrow[l,"e"']
\end{tikzcd}}
Or en dim. $n$ on a \inserted{formule de} changement de base, et \inserted{comme $e'^*(W)\simeq\underline O_B$,} on est ramené à définir un hom.
\[
R^n q_*(\underline O_B)\longrightarrow\deleted{\unsure{\underline O_S|\otimes_S}}\tau_{B^*}
\]
Or on sait que l'on a un isom canonique
\[
R^nq_*(\underline{\mathcal O}_B)\simeq\textstyle\bigwedge^{n}R^1q_*(\underline O_B)\simeq\bigwedge^{n}t_{B^*}
\]
où $t_{B^*}$ est le faisceau tangent de $B^*$ le long de $e$. On \deleted{\unsure{trouve donc}} \inserted{prend alors} \uline{l'hom identique} $\bigwedge^n t_{B^*}\simeq\tau_{B^*}$.
% archive sheet number: 6
\scanpage{8}
Ceci posé, je dis que l'hom précédent
\[
Rp_*(W)\longrightarrow e_*(\underline{\mathcal O}_S)\otimes\tau_{B^*}[-n]
\]
est un isomorphisme. Il suffit pour ceci de prouver que l'homomorphisme induit \deleted{\unsure{sur les|pour les}} en tout point $\unsure{x|b|\deleted{x}b}\in B^*$ par $\mathbb{L} i_x^*$ est un isomorphisme (histoire de mapping-cylinders\ldots). En les points $b\in B^*-e\unsure{_1(S)|(S)|_1(s)}$, ce n'est autre que le lemme 2. En un point $\deleted{b\in e}$ $b$ de la forme $e(s)$, on est ramené au cas où $S$ est le spectre d'un corps algébriquement clos, à faire une vérification de compatibilité\ldots.

\uline{Remarques} \quad 1) a) $\Phi_L$ transforme $e_A$ en $1_B$, $e_s$ ($s$ section de $A/S$) en $L_s$ (faisceau inv. alg. $\sim0$ sur $B$ rel $/S$), \inserted{donc} $e_s*F$ en $L_s\otimes\Phi_L(F)$.

b) $\Phi_L$ transforme \deleted{$\omega_A[n]$ en $\unsure{\varepsilon|e}$, $\omega_A[n]$} $L_t$ ($t$ section de $B$ sur $S$) en $\tau[-n]e_{-t}$ i.e $L_t\,\omega[n]$ en $e_{-t}$, \deleted{\unsure{de|donc}} $F\otimes L_t\deleted{\unsure{\omega[n]}}$ en $\left\{\Phi_L(F)*e_{-t}\right\}\deleted{\unsure{\omega[n]}}\unsure{e|2|,|.}$

c) $\varepsilon_{\unsure{B}}(\Phi_L(F))=\chi_A(F)$, $\chi_B(\Phi_L(F))=(-1)^n\varepsilon_A(F)$
% slanted
\margin{$\varepsilon={}$ augmentation\quad égalité \deleted{\unsure{des|dans} \unsure{fonctions}} $H^0(S,\mathbb Z)$}

2) $\det^*\circ\Phi_L$ définit \deleted{, pour tte section}
\[
\varphi\colon\underline{\mathrm{Pic}}_{A/S}\longrightarrow\underline{\mathrm{Pic}}_{B/S}
\]
qui, pour toute section $\delta$ de $\underline{\mathrm{NS}}_{A/S}$, définit
\[
\underline{\mathrm{Pic}}^\delta_{A/S}\longrightarrow\underline{\mathrm{Pic}}^{-\delta'}_{B/S}\qquad\text{compatible avec}
\]
\[
-(\widetilde{-\delta'})=\widetilde{\delta'}
\]
% archive sheet number: 7
\scanpage{9}
où $\delta'\inserted{\in\underline{\mathrm{NS}}_{B/S}(S)}$ a été déjà caractérisé par ailleurs, par
\[
\widetilde{\delta'}\,\widetilde\delta=-\chi(\delta)\,\mathrm{id}_A\qquad\ldots.
\]

On peut se proposer, pour une section $\lambda$ de $\underline{\mathrm{Pic}}_{A/S}$, définissant le faisceau inv. sur $A$, de ``calculer'' la classe ds $K^{\cdot}(B)$ de $\varphi_L(\deleted{\unsure{\delta}}\mathrm{cl}(L))$ en termes de \deleted{\unsure{la se}} \unsure{de|la} section $\varphi(\lambda)^{\unsure{\,\cdot\,|*|\vphantom{*}}}$ de $\mathrm{Pic}_{B/S}$, définissant un faisceau inversible $L'^{-1}$ sur $B$ (de classe $-\delta'$ dans $\underline{\mathrm{NS}}.(B/S)$ si $\lambda$ est de classe $\delta$ dans $\mathrm{NS}(A/S)$). \unsure{Heuristiqmt|Heuristiquement}, on \deleted{\unsure{aura}} \inserted{devrait avoir} (pour $\chi(\delta)\ne0$ \unsure{partt|partout} sur $S$)
% slanted
\margin{indépendant des résultats de Mumford}
\[
\text{``}\quad\varphi_L(\mathrm{cl}(L))=\chi\deleted{\unsure{\otimes}}L'^{-\frac1\chi}\quad\text{''}
\qquad\text{où }\chi=\chi(\delta)\in H^0(S,\mathbb Z).
\]
Au niveau du $\mathrm{ch}_B$, ceci s'écrit
\margin{$D'\in\mathrm{Gr}^1$ défini par $L'$}
\[
\begin{aligned}
\mathrm{ch}\,\varphi_L(\deleted{\unsure{\mathrm{cl}}}\unsure{\lambda|L|\mathrm{cl}(L)})&=\chi\exp\frac{-D'}{\chi}=\chi-D'+\frac1{2\chi}D'^{2}-\frac1{3!\chi^2}D'^{3}\ldots\\
&\quad\inserted{+(-1)^n\frac1{n!\chi^{n-1}}D'^{n}},
\end{aligned}
\]
et au niveau des classes de Chern elles-mêmes
\[
\begin{aligned}
c(\varphi_L(\unsure{\lambda|L|\mathrm{cl}(L)}))&=\left(1-\frac{D'}{\chi}\right)^\chi\\
&=1-D'+\frac{\chi(\chi-1)}{2\chi}D'^{2}-\frac{\chi(\chi-1)(\chi-2)}{3!\chi^2}D'^{3}+\ldots\\
&\quad+(-1)^n\frac{\chi(\chi-1)\ldots(\chi-n\unsure{+|-}1)}{n!\chi^{n-1}}D'^{n}.
\end{aligned}
\]

\margin{?}
Ces formules sont-elles bien vraies mod torsion ? Elles sont vraies au moins sur un corps alg clos, \uline{mod \uline{équivalence} cohomologique} $l$-adique (tout $l\ne\mathrm{car}.k$). Ceci donne des \margin{!} propriétés de divisibilité intéressantes de $\delta'$ savoir $\deleted{\unsure{\delta'}}\deleted{\unsure{\chi}}\frac{(\chi-1)\ldots(\chi-k+1)}{k!\,\chi^{k-2}}\,\delta'^{k}$ est ``entier'' i.e provient d'un cycle à coeff. entiers
% archive sheet number: 8
\scanpage{10}
3) On vérifie \inserted{par R.R. pour $A\times A\xrightarrow{\pi_A} A$} que l'on a sur $A$
\[
\mathrm{ch}(x*y)=\mathrm{ch}(x)*\mathrm{ch}(y)\operatorname{Todd}(\underline t_A)\tag*{$(*)$}
\]
(donc on a
\[
\mathrm{ch}(x*y)=\mathrm{ch}(x)*\mathrm{ch}\,y
\]
si $\operatorname{Todd}t_{\unsure{A|B}}=1$, p.ex. si $t_A$ est trivial, p.ex. $S$ local).

Considérons d'autre part l'élémt \inserted{$D$} de $\mathrm{Gr}^1(A\times A)$ défini par $L$, et considérons le $\exp D$, d'où
\[
\psi_D=\varphi_{\exp D}\colon\mathrm{Gr}^{\cdot}(A)_{\mathbb Q}\unsure{\mathrel{-}|\to}\mathrm{Gr}^{\cdot}(B)_{\mathbb Q}.
\]

\deleted{\deleted{On voit} On a alors, grâce à R\kern.05em R pour $A\times_S B\unsure{\mathrel{-}|\to}B$ et \deleted{\unsure{cor. précdt} \unsure{cor. de}} prop. 1
\[
\psi_D(xy)=\psi_D(x)\psi_D(y).
\]
On veut utiliser prop. 2 pour calculer $\psi_D(a*\unsure{b|y})$. On a, si
\[
a=\mathrm{ch}\,x,\quad b=\mathrm{ch}\,y
\]
\[
\deleted{\unsure{x}}\ x*y=\mathrm{ch}\,a*\mathrm{ch}\,b=\mathrm{ch}(a*b)T_{\deleted{\unsure{x}}}^{-1}\qquad (T=\operatorname{Todd}(t_A))
\]
d'où
\[
\psi_D(x*y)=\psi_D\bigl(\mathrm{ch}(a*b)T^{-1}\bigr)
\]
}

\deleted{Si} On a, par R\kern.05em R pour $A\times_S B\longrightarrow B$
\[
\mathrm{ch}(\varphi_L(x))=\psi_D(\mathrm{ch}(x))T\qquad (T=\operatorname{Todd}t_A)
\]
\deleted{\unsure{On}} \deleted{\unsure{Posons}} \deleted{Supposons $a=\mathrm{ch}\,x$}

Posons $a=\mathrm{ch}\,x$, $b=\mathrm{ch}\,y$, et transformons $\varphi_L(x*y)=\varphi_L(x)\varphi_L(y)$ (prop 1) par $\mathrm{ch}$, on trouve
\[
\psi_D(\mathrm{ch}(x*y))T=\psi_D(\mathrm{ch}(x))T\,\psi_D(\mathrm{ch}(y))T
\]
\deleted{d'où} d'où, \deleted{\unsure{cf}} grâce à $(*)$
\[
\psi_D(\deleted{\unsure{\mathrm{ch}\,x*}}\mathrm{ch}\,x*\mathrm{ch}\,y)T^{2}\deleted{\unsure{T|T^2}}=\psi_D(\mathrm{ch}\,x)\psi_D(\mathrm{ch}\,y)T^2.
\]
i.e.
\[
\boxed{\psi_D(\unsure{a|\mathrm{ch}}*b)=\psi_D(a)\psi_D(b)}
\]
% archive sheet number: 9
\scanpage{11}
De même, on aura (prop. 2, cor)
\[
\varphi_L(xy)=(\varphi_L(x)*\varphi_L(y))\deleted{\unsure{\mathrm{cl}}}\cdot(-1)^n\omega_A
\]
d'où en appliquant $\mathrm{ch}$ et $({*})$
\[
\psi_D(\mathrm{ch}(xy))=\deleted{\unsure{\psi_D(}}\psi_D(\mathrm{ch}\,x)\,\psi_D(\mathrm{ch}\,y))T(-1)^n\deleted{\exp}\mathrm{ch}\,\omega_A
\]
donc comme $\mathrm{ch}(xy)=\mathrm{ch}\,x\ \mathrm{ch}\,y$
\[
\boxed{\psi_D(ab)=\psi_D(a)\,\psi_D(b)\,(-1)^nT\,\mathrm{ch}\,\omega_A}
\quad\left\{\begin{aligned}
T&=\operatorname{Todd}(t_A)\\
\omega_A&=\bigwedge^n t_A^\vee
\end{aligned}\right.
\]
% slanted
\margin{les formules sont valables dès que $a,b\in\mathrm{Gr}^{\cdot}(A)_{\mathbb Q}$.}

Ce sont les formules \deleted{\unsure{affirmées|différentielles}} \inserted{\unsure{démontrées}} (\unsure{resp|resp.} affirmées) par Kleiman dans ``10 Exposés\ldots'', sur un corps alg clos et en travaillant avec une théorie \unsure{cohomologique|chomologique} plutôt qu'avec $\mathrm{Gr}_{\mathbb Q}$, -- à cela près que Kleiman\deleted{\unsure{n}} n'a pas tenu compte du facteur $(-1)^n$ !

4) Il résulte du Th. de dualité pour $A\times_S B\longrightarrow B$ que l'on a un isom fonctoriel
\[
\Phi_L(\check F\otimes\omega_A[n])\simeq\Phi_L(F)^\vee
\]
i.e.
\[
\boxed{\Phi_L(F)^\vee\simeq\Phi_L(\check F\mathbin{\cdot}\omega_A[n]}
\]
Du point de \unsure{vue|vu} virtuel, ceci s'écrit
\[
\varphi_L(x)^\vee=(-1)^n\omega_A\,\varphi_L(\check x)
\]

5) Bien sûr, si $u$ est un autom de $A$ et $u'$ l'autom contragrédient de $B$, on a par transport de structure un isom canonique
\[
\Phi_L(u.F)\simeq u'.\Phi_L(F)
\]
En particulier, $\Phi_L$ commute à la symétrie $s_A$. Kif-kif pour $\varphi_L$.
% archive sheet number: 10
